(1) 电子稳态运动方程为
\[
0=-e(\boldsymbol E+\boldsymbol v\times\boldsymbol B)-m\boldsymbol v/\tau.
\]
令 $h=eB\tau/m>0,\ \eta=e\tau/m>0$，由于 $\boldsymbol v\times\boldsymbol B=(Bv_y,-Bv_x,0)$，得到
\[
v_x=-\eta E_x-hv_y,\qquad v_y=-\eta E_y+hv_x.
\]
联立解得
\[
v_x=-\frac{\eta(E_x-hE_y)}{1+h^2},\qquad
v_y=-\frac{\eta(hE_x+E_y)}{1+h^2}.
\]
电流密度 $\boldsymbol J=-ne\boldsymbol v$，因此
\[
\binom{J_x}{J_y}=\frac{\sigma_0}{1+h^2}
\begin{pmatrix}1&-h\\h&1\end{pmatrix}
\binom{E_x}{E_y},\quad \sigma_0=ne^2\tau/m.
\]
非对角元素符号来自电子电荷和 $B$ 的固定方向，不能凭习惯调换。

(2) 条带侧边绝缘意味着 $J_y=0$，于是 $hE_x+E_y=0$，即
\[
E_y=-hE_x=-\frac{J_xB}{ne},\qquad J_x=\sigma_0E_x.
\]
若从 $y=0$ 到 $y=w$ 定义电压差 $\Delta\phi=\phi(w)-\phi(0)$，因为 $\boldsymbol E=-\nabla\phi$，
\[
\Delta\phi=-E_yw=\frac{J_xBw}{ne}>0
\]
当 $J_x>0$。这与电子向 $-x$ 漂移相容，电子霍尔系数为 $E_y/(J_xB)=-1/(ne)$。

(3) 磁力始终垂直漂移速度不做功，因此电场供给的功率密度恰好被弛豫碰撞耗散：
\[
p_J=\boldsymbol J\cdot\boldsymbol E=J_xE_x+J_yE_y
=\sigma_0E_x^2.
\]
在侧边开路约束下纵向电阻率 $E_x/J_x=1/\sigma_0$ 与磁场无关，即这个单载流子、各向同性、恒定 $\tau$ 的理想 Drude 模型不存在纵向磁阻；这不是所有真实金属的普适定律。

(4) 如果另一设备强制 $E_y=0$，则
\[
J_x=\frac{\sigma_0}{1+h^2}E_x,\qquad
J_y=\frac{h\sigma_0}{1+h^2}E_x\ne0.
\]
此时侧向载流子可以流动，而纵向电导被因子 $1/(1+h^2)$ 抑制。两种边界条件不同，不能直接把固定 $E_y=0$ 情况的张量分量当作霍尔条带实际测得的纵向电阻率。
